
2.8.
Soit $\underline{L}^{\bullet}$ un complexe parfait sur $X$. Il existe un entier $k$ tel que $\underline{M}^{\bullet}=\mathbb R F_k(\underline{L}^{\bullet})$ soit un complexe $\mathcal C_0$-parfait, et que $\underline{L}^{\bullet}$ soit isomorphe à $\Projz\bigl(\mathbb R F_k(\underline{L}^{\bullet})\bigr)$. Si on applique 2.7 au complexe $\underline{M}^{\bullet}$, la suite des triangles distingués construite donne en appliquant le foncteur $\Projz$ une suite de triangles distingués de $D_{\mathrm{parf}}(\mathcal A)$ (il est clair que $\Projz$ transforme complexe $\mathcal C_0$-parfait en complexe parfait). Par ailleurs, on a :
\[
  \Projz\bigl(\underline{K}^{\bullet}\otimes_{\underline{O}_S}\underline{B}(d)\bigr)
  = f^*(\underline{K}^{\bullet})(d)
\]
(notons que $f^*=\mathbb L f^*$ d'après la platitude de $f$) pour $\underline{K}^{\bullet}$ dans $D^+(\mathcal C')$.

On obtient donc dans $K^{\bullet}(X)$ l'égalité :
\[
  \cl(\underline{L}^{\bullet})
  =\sum_i\cl\bigl(f^*(\underline{K}_i^{\bullet})\bigr).\xi^{d_i}
  =\sum_i f^*\bigl(\cl(\underline{K}_i^{\bullet})\bigr).\xi^{d_i},
\]
où les $\underline{K}_i^{\bullet}$ sont les complexes parfaits sur $S$ définis dans 2.7. Il en résulte que $K^{\bullet}(X)$ est engendré comme $K^{\bullet}(S)$-module par les puissances de $\xi$.

Pour voir que $1,\xi,\ldots,\xi^{r+1}$ sont linéairement dépendants, on considère l'homomorphisme~: $K^{\bullet}_{\naif}(X)\to K^{\bullet}_{\mathrm{parf}}(X)$ qui à la classe d'un faisceau localement libre de type fini $\underline{L}$ associe la classe du complexe égal à $\underline{L}$ en degré $0$, et nul pour les autres degrés. L'équation (1.9.1) donne alors une relation de dépendance dans $K^{\bullet}_{\mathrm{parf}}(X)$.

Pour voir que $1,\xi,\ldots,\xi^r$ sont linéairement indépendants, on peut utiliser le morphisme $K^{\bullet}(S)$-linéaire $f_*:K^{\bullet}(X)\to K^{\bullet}(S)$ (IV 2.11.1). On a :
\[
  f_*(1)=1,
  \quad\text{car } f_*(\underline{O}_X)=\underline{O}_S,\text{ et }R^i(f_*(\underline{O}_X))=0\text{ pour }i\geq1;
\]
\[
  f_*(\xi^n)=0\quad\text{pour }-r\leq n<0,
  \quad\text{car }R^i\bigl(f_*(\underline{O}_X(n))\bigr)=0\text{ pour }i\geq0.
\]
Si on a une relation du type :
\[
  f^*(x_0)+f^*(x_1)\xi+\ldots+f^*(x_r)\xi^r=0,
\]
on la multiplie par $\xi^{-r}$, et on applique $f_*$; on trouve alors $x_r=0$; on 
